\section{Where Ymir fits best}\label{sec:fit}

Concrete scenarios where Ymir's built features give it an edge over the alternatives above, with
the honest catch for each. Anything that relies on unbuilt work is marked planned. Ymir has no
tagged release yet, so every scenario currently means building from source.

\begin{rowlist}
  \item[1. Privacy-sensitive work on hardware you own \status{Built}]
  Legal drafts, medical notes or proprietary source code. Why Ymir: prompts, answers and synced context move only between your linked devices over Noise-encrypted streams, are encrypted at rest, and are never logged as text; there is no endpoint to leak. Where others fall short: hosted APIs, Microsoft's hybrid Copilot and Apple PCC all send some requests to vendor servers (Apple's are attested and verifiable, but still Apple's); Ollama or llama-swap keep data local but need a proxy and keys before another machine can reach them.

  \item[2. A home lab of mismatched machines \status{Built}]
  An NVIDIA desktop, an AMD box, an Apple Silicon Mac and a CPU-only server. Why Ymir: each node runs the engine that suits it (llama.cpp with CUDA or Vulkan, MLX-LM on Apple Silicon, vLLM on Linux NVIDIA), advertises what it can run, and joins one catalog with load balancing. Where others fall short: llama-swap and Ollama manage one machine each; exo pools devices on one network with an MLX focus and centres on sharding one model; GPUStack and data-center stacks expect clusters, Kubernetes or a control plane.

  \item[3. Your phone on your home GPUs, with no ports or VPN to run \status{Built}]
  Why Ymir: Bragi on Android or iOS pairs with your identity by QR code and reaches inference nodes over Weirdnet, using its relay and direct paths, so nothing is port-forwarded and no VPN client is needed. Where others fall short: exposing Ollama means opening an unauthenticated API or adding a proxy; LM Studio LM Link avoids public ports via a Tailscale mesh but links devices running LM Studio and is in preview; Firebase's Gemini Nano hybrid falls back to Google's cloud rather than your machines.

  \item[4. Long answers and agent runs that survive a node dying \status{Built (job hand-off planned)}]
  Why Ymir: if a model node or an inference node fails mid-answer, the reply continues on another node from the exact offset with no repeated text, up to two resumes per level; agent tasks for Bragi run on the pool's models. Where others fall short: Ollama and llama-swap have no failover; LocalAI federation moves a failed request to the next member; NVIDIA Dynamo offers comparable request migration but is built for Kubernetes GPU fleets. Planned: a batch or agent job whose submitting node goes down is picked up by another node from a synced job record.

  \item[5. Heavy batch work instead of per-token bills \status{Built}]
  Summarising a document archive or labelling a dataset overnight. Why Ymir: \code{ymirctl batch} fans prompts across every node at batch priority with retries, while device fairness caps keep a phone's interactive chat responsive. There is no per-token fee; the cost is hardware and power you already have. Where others fall short: hosted APIs meter usage by tokens, and hybrid services such as GitHub Copilot Fusion still send the work they route to the cloud through a paid plan; single-machine servers cannot spread one job over several machines.

  \item[6. Local tools on pool models \status{Built}]
  Jan, editors, scripts and MCP clients such as Cursor or VS Code. Why Ymir: the OpenAI-compatible bridge on 127.0.0.1 and \code{ymirctl mcp} give every local tool the whole pool's catalog, with a local token and no network listener; Bragi's VS Code and Cursor extension and \code{bragi mcp} reach the pool the same local way. Where others fall short: Ollama and llama-swap serve only the models on that one machine; LM Link offers similar local access but only to LM Studio devices.

  \item[7. Mostly offline or locked-down networks \status{Built, release pending}]
  Why Ymir: release files are designed to bundle a pinned llama.cpp so setup needs no extra install, models download only with consent, nodes find each other on the LAN by mDNS, and the relay used for bootstrapping can be pointed at your own. Where others fall short: GitHub notes Copilot Fusion makes routing decisions in the cloud even for local models; hosted APIs need the internet. Catch: Ymir bootstraps from the public relay by default, and the bundled release files are not published yet.

  \item[8. A small studio or office sharing a few machines \status{Fits with limits}]
  Why Ymir can fit: several GPU boxes become one endpoint with fair sharing between calling devices. Catch: Ymir has a single owner by design, with no user accounts, roles or quotas, so it only fits when every device belongs to the same identity. Teams that need per-user access control are better served by GPUStack, which has multi-tenancy and RBAC.
\end{rowlist}

Where Ymir is the wrong choice: serving many users or customers, running frontier-scale models it
cannot fit on your hardware, or squeezing maximum throughput out of a GPU fleet. Hosted APIs, the
vendor hybrids and data-center stacks such as Dynamo and llm-d are built for those jobs.
